% GENERATED FILE. Edit canonical lessons, then run npm run generate:books.
% canonical-source-hash: fnv1a64:e6aed79fd604f767
% canonical-lessons: KA-C119-ole, KA-C119-menabatti, KA-C119-benkipettige, KA-C119-bombe, KA-C119-cendu

\chapter{Stove, Candle, Matchbox, Doll, Ball}
\label{ch:ka-stove-candle-matchbox-doll-ball}
\hlchaptermodality{119}

\begin{chapteropening}
\textbf{By the end of this chapter:} \emph{I can say a stove, a candle, a matchbox, a doll and a ball.}

Complete the last lesson of chapter 119: i can say a stove, a candle, a matchbox, a doll and a ball.
\end{chapteropening}

\section[\texorpdfstring{\kn{ಒಲೆ} (ole)}{ಒಲೆ (ole)}]{\kn{ಒಲೆ} (ole) — a stove}

\label{lesson:KA-C119-ole}



\begin{quote}
Before the new one: say the Kannada for a pen, then the Kannada for paper; a letter.
\end{quote}

\subsection*{You'll want to know: \kn{ಒಲೆ}}
\textbf{\kn{ಒಲೆ}} — \emph{ole} — \textquotedblleft{}a stove\textquotedblright{}.

\begin{grammarlens}[title={What you've built}]
One more word for this chapter.
\end{grammarlens}

\subsection*{Your turn}
\begin{itemize}
\raggedright
  \item \emph{Say it:} \emph{ole}
  \item \emph{Say it:} \emph{ole}, once more
  \item \emph{Say it:} say \emph{ole}
  \item \emph{Recall:} say the Kannada for a pen, then the Kannada for paper; a letter, then say \emph{ole} again
\end{itemize}

\begin{tcolorbox}[breakable,colback=teal!4,colframe=teal!35!black,title={Before you move on}]
What does \kn{ಒಲೆ} mean? (\textquotedblleft{}A stove\textquotedblright{}.) Say it once more.
\end{tcolorbox}
\section[\texorpdfstring{\kn{ಮೇಣಬತ್ತಿ} (mēṇabatti)}{ಮೇಣಬತ್ತಿ (mēṇabatti)}]{\kn{ಮೇಣಬತ್ತಿ} (mēṇabatti) — a candle}

\label{lesson:KA-C119-menabatti}



\begin{quote}
Before the new one: say the Kannada for paper; a letter, then the Kannada for a stove.

Greet through the day: \textbf{\kn{ಶುಭೋದಯ}}, \textbf{\kn{ಶುಭ} \kn{ಮಧ್ಯಾಹ್ನ}}, \textbf{\kn{ಶುಭ} \kn{ಸಂಜೆ}}. Then name the afternoon, \textbf{\kn{ಅಪರಾಹ್ನ}}.
\end{quote}

\subsection*{You'll want to know: \kn{ಮೇಣಬತ್ತಿ}}
\textbf{\kn{ಮೇಣಬತ್ತಿ}} — \emph{mēṇabatti} — \textquotedblleft{}a candle\textquotedblright{}.

\begin{grammarlens}[title={What you've built}]
One more word for this chapter.
\end{grammarlens}

\subsection*{Your turn}
\begin{itemize}
\raggedright
  \item \emph{Say it:} \emph{mēṇabatti}
  \item \emph{Say it:} \emph{mēṇabatti}, once more
  \item \emph{Say it:} say \emph{mēṇabatti}
  \item \emph{Recall:} say the Kannada for paper; a letter, then the Kannada for a stove, then say \emph{mēṇabatti} again
\end{itemize}

\begin{tcolorbox}[breakable,colback=teal!4,colframe=teal!35!black,title={Before you move on}]
What does \kn{ಮೇಣಬತ್ತಿ} mean? (\textquotedblleft{}A candle\textquotedblright{}.) Say it once more.
\end{tcolorbox}
\section[\texorpdfstring{\kn{ಬೆಂಕಿಪೆಟ್ಟಿಗೆ} (beṅkipeṭṭige)}{ಬೆಂಕಿಪೆಟ್ಟಿಗೆ (beṅkipeṭṭige)}]{\kn{ಬೆಂಕಿಪೆಟ್ಟಿಗೆ} (beṅkipeṭṭige) — a matchbox}

\label{lesson:KA-C119-benkipettige}



\begin{quote}
Before the new one: say the Kannada for a stove, then the Kannada for a candle.
\end{quote}

\subsection*{You'll want to know: \kn{ಬೆಂಕಿಪೆಟ್ಟಿಗೆ}}
\textbf{\kn{ಬೆಂಕಿಪೆಟ್ಟಿಗೆ}} — \emph{beṅkipeṭṭige} — \textquotedblleft{}a matchbox\textquotedblright{}.

\begin{grammarlens}[title={What you've built}]
One more word for this chapter.
\end{grammarlens}

\subsection*{Your turn}
\begin{itemize}
\raggedright
  \item \emph{Say it:} \emph{beṅkipeṭṭige}
  \item \emph{Say it:} \emph{beṅkipeṭṭige}, once more
  \item \emph{Say it:} say \emph{beṅkipeṭṭige}
  \item \emph{Recall:} say the Kannada for a stove, then the Kannada for a candle, then say \emph{beṅkipeṭṭige} again
\end{itemize}

\begin{tcolorbox}[breakable,colback=teal!4,colframe=teal!35!black,title={Before you move on}]
What does \kn{ಬೆಂಕಿಪೆಟ್ಟಿಗೆ} mean? (\textquotedblleft{}A matchbox\textquotedblright{}.) Say it once more.
\end{tcolorbox}
\section[\texorpdfstring{\kn{ಬೊಂಬೆ} (bombe)}{ಬೊಂಬೆ (bombe)}]{\kn{ಬೊಂಬೆ} (bombe) — a doll}

\label{lesson:KA-C119-bombe}



\begin{quote}
Before the new one: say the Kannada for a candle, then the Kannada for a matchbox.
\end{quote}

\subsection*{You'll want to know: \kn{ಬೊಂಬೆ}}
\textbf{\kn{ಬೊಂಬೆ}} — \emph{bombe} — \textquotedblleft{}a doll\textquotedblright{}.

\begin{grammarlens}[title={What you've built}]
One more word for this chapter.
\end{grammarlens}

\subsection*{Your turn}
\begin{itemize}
\raggedright
  \item \emph{Say it:} \emph{bombe}
  \item \emph{Say it:} \emph{bombe}, once more
  \item \emph{Say it:} say \emph{bombe}
  \item \emph{Recall:} say the Kannada for a candle, then the Kannada for a matchbox, then say \emph{bombe} again
\end{itemize}

\begin{tcolorbox}[breakable,colback=teal!4,colframe=teal!35!black,title={Before you move on}]
What does \kn{ಬೊಂಬೆ} mean? (\textquotedblleft{}A doll\textquotedblright{}.) Say it once more.
\end{tcolorbox}
\section[\texorpdfstring{\kn{ಚೆಂಡು} (ceṇḍu)}{ಚೆಂಡು (ceṇḍu)}]{\kn{ಚೆಂಡು} (ceṇḍu) — a ball}

\label{lesson:KA-C119-cendu}



\begin{quote}
Before the new one: say the Kannada for a matchbox, then the Kannada for a doll.
\end{quote}

\subsection*{You'll want to know: \kn{ಚೆಂಡು}}
\textbf{\kn{ಚೆಂಡು}} — \emph{ceṇḍu} — \textquotedblleft{}a ball\textquotedblright{}.

\begin{grammarlens}[title={What you've built}]
One more word for this chapter.
\end{grammarlens}

\subsection*{Your turn}
\begin{itemize}
\raggedright
  \item \emph{Say it:} \emph{ceṇḍu}
  \item \emph{Say it:} \emph{ceṇḍu}, once more
  \item \emph{Say it:} say \emph{ceṇḍu}
  \item \emph{Recall:} say the Kannada for a matchbox, then the Kannada for a doll, then say \emph{ceṇḍu} again
\end{itemize}

\begin{tcolorbox}[breakable,colback=teal!4,colframe=teal!35!black,title={Before you move on}]
What does \kn{ಚೆಂಡು} mean? (\textquotedblleft{}A ball\textquotedblright{}.) Say it once more.
\end{tcolorbox}
